% Clase
\documentclass[11pt,a4paper,spanish,twoside]{book}

% Órdenes auxiliares
% Español
\usepackage[spanish]{babel}
\usepackage[utf8]{inputenc}
\usepackage[T1]{fontenc}
\usepackage{lmodern}

% Imágenes
\usepackage[pdftex]{graphicx}
\usepackage{latexsym}
\usepackage{fancybox}

% Ruta para las imágenes
\graphicspath{{img/}}

% Rotaciones
\usepackage[twoside]{rotating}

% Multirow
\usepackage{multirow}

% Referencias
\usepackage[spanish]{varioref}
\usepackage[pdftex,colorlinks=true,linkcolor=black]{hyperref}

% Colores
\usepackage{color}
\usepackage{colortbl}
\definecolor{rojo}{rgb}{1.0,0.0,0.0}

% Párrafos
\setlength{\parskip}{6pt}

% Code for creating empty pages
% No headers on empty pages before new chapter
\makeatletter
\def\cleardoublepage{\clearpage\if@twoside \ifodd\c@page\else
    \hbox{}
    \thispagestyle{empty}
    \newpage
    \if@twocolumn\hbox{}\newpage\fi\fi\fi}
\makeatother \clearpage{\pagestyle{empty}\cleardoublepage}

%% Entorno Listings para código fuente
\usepackage{listingsutf8}[2007/11/11]

\lstset{
  frame=Ltb, %forma del cuadro
  framerule=0pt, %ancho del borde
  aboveskip=0.5cm, %separación de los números de línea
  framexleftmargin=0.4cm, %margen externo izquierdo
  framesep=0pt,
  rulesep=.4pt,
  rulesepcolor=\color{black},
  % 
  stringstyle=\ttfamily,
  showstringspaces = false,
  basicstyle=\scriptsize,
  keywordstyle=\bfseries,
  % 
  numbers=left,
  numbersep=15pt,
  numberstyle=\tiny,
  numberfirstline= false,
  %
  inputencoding=utf8/latin1
}

% minimizar fragmentado de listados
\lstnewenvironment{listing}[1][]{
  \lstset{#1}\pagebreak[0]}{\pagebreak[0]
}


% Fancyhdr - Encabezados y pies de página
\usepackage{fancyhdr}
% Márgenes
\headsep=8mm
\footskip=14mm

% Fancy Header Style Options
\pagestyle{fancy}               % Sets fancy header and footer
\fancyfoot{}                    % Delete current footer settings

% Sin mayúsculas en la cabecera
\lhead{\nouppercase{\rightmark}}
\rhead{\nouppercase{\leftmark}}

% Capítulo
\renewcommand{\chaptermark}[1]{ % Lower Case Chapter marker style
  \markboth{\chaptername\ \thechapter.\ #1}{}} 

% Sección
\renewcommand{\sectionmark}[1]{ % Lower case Section marker style
  \markright{\thesection.\ #1}} 

% Página
\fancyhead[LE,RO]{\bfseries\thepage} % Page number (boldface) in left on even
                                     % pages and right on odd pages

% ------------------ Macro para encabezados y pies de página-------------------
%    Uso: \encabezado{pares(pag izquierda)}
% -----------------------------------------------------------------------------
\def\encabezado{
  \fancyhead[RE]{\bfseries\leftmark}     % In the right on even pages
  \fancyhead[LO]{\bfseries\rightmark}  % In the left on odd pages
  \renewcommand{\headrulewidth}{0.5pt} % Width of head rule
}
% -----------------------------------------------------------------------------

% ------------------ Macro para insertar una imagen ---------------------------
%    Uso: \imagen{nombreFichero}{Ancho(cm))}{Etiqueta}{Identificador}
% -----------------------------------------------------------------------------
\usepackage{float}
\usepackage{ifthen}
\def\imagen#1#2#3#4{
  \begin{figure}[H]
    \begin{center}
      \ifthenelse{\equal{#2}{}}
      {\includegraphics{#1}}{\resizebox{#2cm}{!}{\includegraphics{#1}}}
      \ifthenelse{\equal{#3}{}}{}{\caption{#3}}
      \label{#4}
    \end{center}
  \end{figure}
}
% -----------------------------------------------------------------------------


% ------------------ Macro para la portada ------------------------------------
%    Uso: \portada{asignatura}{titulo}{subtítulo}{autor}{fecha}
% -----------------------------------------------------------------------------
\def\portada#1#2#3#4#5{
  \thispagestyle{empty}
  \vspace*{-3.3cm}

  \begin{minipage}[t]{14cm}
    \begin{center}
      \includegraphics[scale=0.25]{logoesi}
  
      \vspace*{1.5cm}
      {\Large \textbf{Universidad de Castilla-La Mancha\\ 
          Escuela Superior de Informática}\\}
    
      \vspace*{1.2cm}
      {\huge \textbf{#1}\\}
    
      \vspace*{1.5cm}
      {\huge #2}\\{\Large #3\\}
    
      \vspace*{1.5cm}
      {\large #4\\}
      \vspace*{1.4cm}
      \large{#5}
    \end{center}
  \end{minipage}

  \newpage
%   \vspace*{1cm}
%   \thispagestyle{empty} 
%   \newpage
}
% -----------------------------------------------------------------------------


% ------------------ Macro para la licencia -----------------------------------
%    Uso: \portada{asignatura}{titulo}{subtítulo}{autor}{fecha}
% -----------------------------------------------------------------------------
\def\licencia#1{
  \thispagestyle{empty}  % Suprime la numeración de esta página
  \vspace*{16.5cm}
  \begin{small}
    \copyright~ #1. Se permite la copia, distribución y/o 
    modificación de este documento bajo los términos de la licencia de
    documentación libre GNU, versión 1.1 o cualquier versión posterior publicada
    por la {\em Free Software Foundation}, sin secciones invariantes. Puede
    consultar esta licencia en http://www.gnu.org. \\[0.2cm]
    Este documento fue compuesto con \LaTeX{}. 
  \end{small}
  \newpage
%   \thispagestyle{empty}
%   \vspace*{0cm}
%   \newpage
}
% -----------------------------------------------------------------------------




% Encabezado y pie de página
\encabezado

\usepackage{enumerate}

\begin{document}

% Silabación extra
\hyphenation{
a-sig-na-tu-ras
au-to-ma-ti-za-rá
ca-tá-lo-go
ca-rre-ra
diag-nos-tico
in-fe-rior
man-te-ni-mien-to
per-so-nal
pro-por-cio-na-rá
pu-bli-ca-da
re-qui-si-tos
res-pecto
u-su-a-rios
vi-lla-rre-al
}


% Portada
\portada{Planificación y Gestión de\\Sistemas de Información}
{Trabajo 5}{Estimación del software}
{Sergio de la Rubia García-Carpintero\\Miguel Millán Sánchez-Grande\\
  Luis Muñoz Villarreal\\Alicia Serrano Sánchez\\
  Juan Miguel Torres Triviño}{24 de Mayo de 2010}

% Licencia
\licencia{Sergio de la Rubia García-Carpintero, Miguel Millán Sánchez-Grande,
  Luis Muñoz Villarreal, Alicia Serrano Sánchez, Juan Miguel Torres Triviño}

\chapter*{Ficha de trabajo}
\begin{description}
\item[Código] T5
\item[Fecha] 24 de Mayo de 2010
\item[Título] Estimación del software
\end{description}

\begin{table}[!ht]
  \centering
  \begin{tabular}{lp{5cm}c}
    \multicolumn{3}{l}{\Large \textbf{Equipo} G4} \\ \\
    \multicolumn{1}{c}{\emph{Apellidos y nombre}} & 
    \multicolumn{1}{c}{\emph{Firma}} & \emph{Puntos} \\
    \hline \\
    de la Rubia García-Carpintero, Sergio & & 8 \\ \\
    Millán Sánchez-Grande, Miguel         & & 8 \\ \\
    Muñoz Villarreal, Luis                & & 8 \\ \\
    Serrano Sánchez, Alicia               & & 8 \\ \\
    Torres Triviño, Juan Miguel           & & 8 \\ \\
    \hline
  \end{tabular}
\end{table}

% Índices
\tableofcontents
\listoftables
%\listoffigures

%% INICIO DEL DOCUMENTO %%%%%%%%%%%%%%%%%%%%%%%%%%%%%%%%%%%%%%%%%%%%%%%%%
\chapter*{Introducción}
En este documento se realiza la estimación del esfuerzo del software para
realizar guías docentes. La estimación del software es la base que servirá
para poder realizar una estimación de costo en la etapa posterior.

Los principales módulos que se identifican en el software que se
va a desarrollar son \emph{gestión de usuarios}, \emph{gestión de cursos} e
\emph{interfaz}. Para obtener los puntos función del software sin ajustar se
utiliza el método \emph{Puntos Función} definido por Allan Albrecht
\cite{albrecht} y empleando la metodología de medición de la
\emph{International Function Point Users Group}

En cada módulo del sistema software se identifican entradas externas, salidas
externas, consultas externas, archivos lógicos internos y archivos lógicos de
interfaz externos para estimar el tamaño de cada uno de ellos.

\section*{Visión general}
El desarrollo del sistema pretende facilitar el control de todos los aspectos
relacionados con la generación automatizada de guías docentes para la
Universidad de Málaga \cite{uma}. El software tiene funcionalidad para
construir una 
guía docente a partir de los datos almacenados en las distintas bases de
datos de la universidad, así como adaptar los planes de estudio antiguos a
los nuevos acordados por el EEES\footnote{El Espacio Europeo de Educación
  Superior es un ámbito de organización educativo iniciado en 1999 con la
  \emph{Declaración de Bolonia} que quiere armonizar los distintos sistemas
  educativos de la Unión Europea. Fuente:
  http://es.wikipedia.org/wiki/Espacio\_Europeo\_de\_Educación\_Superior},
asignar el profesorado a las distintas asignaturas, el establecimiento de los
horarios lectivos y los horarios de exámenes.

\section*{Objetivos}
El objetivo principal es la agilización de la realización de guías
docentes. Este objetivo principal conlleva otros objetivos que también son
importantes para este sistema, como la actualización en tiempo real de los
contenidos de las guías docentes, mejorar la accesibilidad a dichos
contenidos y la reducción en el gasto de la contratación de personal para la
elaboración de estas guías.

\section*{Procedimiento}
Para llevar a cabo esta medición se seguirá la metodología de puntos función
del \emph{IFPUG}, que contiene los siguientes pasos:
\begin{enumerate}
\item Identificar la frontera de la aplicación. \\
  \emph{Capítulo \ref{cap:mod}: \nameref{cap:mod}}
\item Identificar los puntos funcionales. \\
  \emph{Capítulo \ref{cap:mod}: \nameref{cap:mod}}
\item Evaluar la complejidad. \\
  \emph{Capítulo \ref{cap:com}: \nameref{cap:com}}
\item Calcular los PF sin ajustar. \\
  \emph{Capítulo \ref{cap:pun}: \nameref{cap:pun}} 
\item Evaluar los atributos de ajuste. \\ 
  \emph{Capítulo \ref{cap:fac}: \nameref{cap:fac}} 
\item Calcular el valor de ajuste. \\
  \emph{Capítulo \ref{cap:fac}: \nameref{cap:fac}} 
\item Calcular el valor final de los PF. \\
  \emph{Capítulo \ref{cap:aju}: \nameref{cap:aju}}
\end{enumerate}

\chapter{Módulos del software} \label{cap:mod}
\section{Gestión de usuarios}
\subsection{Entradas externas}
\subsubsection{Identificación del usuario}
Cada usuario debe identificarse mediante un nombre y una contraseña.

\subsubsection{Preferencia de horarios} 
Cada profesor introduce sus preferencias de horarios de trabajo para que el
sistema lo tenga en cuenta a la hora de la realización del horario oficial
del curso.

\subsection{Salidas externas}
\subsubsection{Datos personales}
Muestra los datos personales de un determinado profesor, incluyendo
nombre, correo electrónico, teléfono, publicaciones, currículum, \dots

\subsubsection{Horario de un profesor}
El sistema muestra el horario que tiene asignado un determinado profesor.

\subsubsection{Tutorías de los profesores} 
Muestra los horarios de tutorías junto con el despacho asignado de los
profesores de cada facultad.
 
\subsubsection{Asignaturas impartidas por cada profesor}
El sistema muestra la información de cada asignatura que comprende:
\begin{itemize}
\item Profesor/es que la imparte/n.
\item Créditos.
\item Contenidos de la asignatura.
\item Asignaturas recomendadas.
\item Página web de la asignatura.
\item Planificación docente.
\item Sistema de evaluación.
\end{itemize}

\subsection{Consultas externas}
\subsubsection{Consultar currículum}
Información del currículum de un determinado profesor.

\subsubsection{Consultar datos de asignatura} 
Información sobre una determinada asignatura (créditos, profesor que la
imparte, planificación, \dots).

\subsubsection{Horarios de asignaturas}
Horarios de las asignaturas de cada curso. Cada horario
también muestra información de la localización en la que cada asignatura se
imparte. 

\subsection{Archivos lógicos de interfaz externos}
\subsubsection{Datos de los profesores}
El sistema tiene una base de datos donde se guarda toda la información con
respecto a los profesores (nombre, apellidos, teléfono, correo electrónico,
asignaturas que imparte, \dots).

\subsubsection{Datos de las asignaturas}
En la base de datos de la universidad se tienen los datos de las asignaturas,
tales como nombre, créditos, guión de la asignatura, contenidos, evaluación,
\dots 


\section{Gestión de cursos}
\subsection{Entradas externas}
\subsubsection{Añadir curso}
EFntrada destinada a la inclusión de los distintos cursos correspondientes al
plan de estudios en el sistema.

\subsubsection{Añadir asignaturas}
Entrada generada para la inserción de las diferentes asignaturas asociadas a
cada curso.

\subsubsection{Añadir aulas}
Entrada destinada a la introducción de las aulas disponibles destinadas al
correspondiente plan de estudio.

\subsubsection{Añadir profesores}
Entrada generada para incluir los profesores encargados de la docencia de
cada asignatura.

\subsection{Salidas externas}
\subsubsection{Relación cursos/asignaturas}
Información sobre las asignaturas que forman cada curso.

\subsubsection{Relación asignaturas/profesores}
Información de los profesores que imparten la docencia de cada asignatura.

\subsubsection{Relación asignaturas/aulas}
Información sobre la relación entre las asignaturas y las aulas donde son
impartidas. 

\subsubsection{Cuadrante de horarios}
Tabla de horarios por cursos con la información de cada asignatura, aula y
profesor que la imparte.

\subsection{Consultas externas}
\subsubsection{Consultar relación cursos/asignaturas}
Consulta de la información necesaria para relacionar las tablas de
cursos con las tablas de asignaturas.

\subsubsection{Consultar relación asignaturas/profesores}
Consulta de la información necesaria para relacionar las tablas de
asignaturas con las tablas de profesores.

\subsubsection{Consultar relación cursos/aulas}
Consulta de la información necesaria para relacionar las tablas de
cursos con las tablas de aulas.

\subsubsection{Consultar cuadrante de horarios}
Consulta de la información requerida para realizar los cuadrantes de horarios.

\subsection{Archivos lógicos internos}
\subsubsection{Datos de las asignaturas}
Toda la información relacionada con cada asignatura, temario, créditos \dots

\subsubsection{Datos de los cursos}
Información relacionada con cada curso así como las asignaturas que los forman.

\subsubsection{Datos de las aulas}
Información de las aulas disponibles y las utilizadas para la docencia de las
asignaturas. 

\subsection{Archivos lógicos de interfaz externos}
\subsubsection{Historial de guías docentes}
Datos externos relacionados con las guías docentes de cursos anteriores. 
 
\section{Interfaz}
\subsection{Entradas externas}
\subsubsection{Autentificación de los usuarios}
Recoge los datos necesarios para ver si un usuario entra al sistema,
dependiendo si este es administrador o no.

\subsubsection{Introducción de datos para la inserción/modificación/borrado 
de información de un determinado usuario}
Por su similitud se explican conjuntamente pero equivale a tres
\emph{funciones de usuario} diferentes. El objetivo de estas funciones es
recoger los datos necesarios para la de inserción, modificación o borrado de
la información de un usuario en la aplicación.

\subsection{Salidas externas}
\subsubsection{Mostrar los datos para una búsqueda de información con
  posibilidad de ordenación por diferentes criterios}
El objetivo de este proceso es mostrar un listado de usuarios que cumplen las
condiciones de búsqueda solicitadas y en el orden indicado.

\subsubsection{Mostrar los datos de un determinado usuario}
El objetivo de este proceso es mostrar información detallada referente a un
determinado usuario (profesores).

\subsubsection{Mostrar los datos de un determinado curso}
El objetivo de este proceso es mostrar información detallada referente a un
curso, incluyendo horario, aulas, asignaturas, etc.

\subsection{Consultas externas}
\subsubsection{Búsqueda de información con posibilidad de ordenación por
  diferentes criterios}
El objetivo de este proceso es la introducción de unos criterios para realizar 
una búsqueda detallada tanto de cursos como de usuarios.

\subsubsection{Consulta de los datos de un determinado usuario}
El objetivo de esta función es la consulta simple de la información referente
a un usuario (profesor).

\subsubsection{Consulta de los datos de un determinado curso}
El objetivo de esta función es la consulta simple de la información referente
a un curso.

\subsection{Archivos lógicos internos}
\subsubsection{Importación de datos desde un fichero}
El objetivo de esta función es introducir los datos necesarios para importar 
inserciones, modificaciones o borrados desde un fichero de texto.

\subsubsection{Exportación de los resultados de la búsqueda a un fichero}
El objetivo de esta función es introducir los datos para obtener de la base
de datos de la aplicación el listado de usuarios que cumplen las condiciones de
búsqueda solicitadas y en el orden indicado en un fichero de texto.

\subsection{Archivos lógicos de interfaz externos}
No se ha detectado ningún archivo lógico de interfaz externo.

\chapter{Complejidad de las funciones de los  módulos} \label{cap:com}
\section{Gestión de usuarios}
\subsection{Entradas externas}
\subsubsection{Identificación del usuario}
La complejidad de esta función de usuario es \textbf{baja} porque se hace
referencia a un registro para la información de usuario y el número de datos
elementales es dos. 

\subsubsection{Preferencia de horarios} 
La complejidad de esta función es \textbf{media} porque se hace
referencia a dos registros donde se tienen guardados los horarios y las
asignaturas y el número de datos elementales son que se utilizan son seis. 

\subsection{Salidas externas}
\subsubsection{Datos personales}
La complejidad de esta función es \textbf{media} porque se hace referencia a
un registro para la información de usuario y los datos elementales que se
utilizan son dieciséis.

\subsubsection{Horario de un profesor}
La complejidad de esta función es \textbf{media} porque se hace referencia a
dos registros, horarios y asignaturas; y los datos elementales que
se utilizan son seis.

\subsubsection{Tutorías de los profesores} 
La complejidad de esta función es \textbf{baja} porque se hace referencia a
dos registros, profesores y tutorías; y el número de datos elementales que se
utilizan son cuatro. 

\subsubsection{Asignaturas impartidas por cada profesor}
La complejidad de esta función es \textbf{medio} porque se hace referencia a
dos registros, profesores y asignaturas; y el número de datos elementales que se
utilizan son siete. 

\subsection{Consultas externas}
\subsubsection{Consultar currículum}
La complejidad de esta función es \textbf{baja} porque se hace referencia a
un registro y el número de datos elementales es dos.

\subsubsection{Consultar datos de asignatura} 
La complejidad de esta función es \textbf{baja} porque se hace referencia a
un registro y el número de datos elementales es nueve.

\subsubsection{Horarios de asignaturas}
La complejidad de esta función es \textbf{baja} porque se hace referencia a
tres registros y los datos elementales que se utilizan son cuatro.

\subsection{Archivos lógicos de interfaz externos}
\subsubsection{Datos de los profesores}
La complejidad de esta función es \textbf{baja} porque se hace referencia a
un registro y los datos elementales son siete.

\subsubsection{Datos de las asignaturas}
La complejidad de esta función es \textbf{baja} porque se hace referencia a
un registro y los datos elementales son catorce.

\section{Gestión de cursos}
\subsection{Entradas externas}
\subsubsection{Añadir curso}
La complejidad de esta función es \textbf{baja}, porque tiene un
registro para la información del curso y el número de tipo de datos
elementales es menor que quince.

\subsubsection{Añadir asignaturas}
La complejidad de esta función es \textbf{baja}, porque tiene un
registro para la información de las asignaturas y el número de tipo de datos
elementales es menor que quince.

\subsubsection{Añadir aulas}
La complejidad de esta función es \textbf{baja}, porque tiene un
registro para la información de las aulas y el número de tipo de datos
elementales es menor que quince.

\subsubsection{Añadir profesores}
La complejidad de esta función es \textbf{baja}, porque tiene un
registro para la información de los profesores y el número de tipo de datos
elementales es menor que quince.

\subsection{Salidas externas}
\subsubsection{Relación cursos/asignaturas}
La complejidad de esta función es \textbf{media}, porque usa dos
registros para la información del curso y las asignaturas y el número de tipo
de datos elementales se encuentra entre seis y diecinueve.

\subsubsection{Relación asignaturas/profesores}
La complejidad de esta función es \textbf{media}, porque usa dos
registros para la información de las asignaturas y los profesores y el número
de tipo de datos elementales se encuentra entre seis y diecinueve. 

\subsubsection{Relación asignaturas/aulas}
La complejidad de esta función es \textbf{media}, porque usa dos
registros para la información de las aulas y las asignaturas y el número de
tipo de datos elementales se encuentra entre seis y diecinueve.

\subsubsection{Cuadrante de horarios}
La complejidad de esta función es \textbf{alta}, porque usa todos los
registros de este módulo para mostrar el cuadrante y el número de
tipo de datos elementales es mayor que veinte.

\subsection{Consultas externas}
\subsubsection{Consultar relación cursos/asignaturas}
La complejidad de esta función es \textbf{media}, porque usa dos
registros para la información de las asignaturas y los cursos de este
módulo y el número de tipo de datos elementales se encuentra entre seis y
diecinueve.

\subsubsection{Consultar relación asignaturas/profesores}
La complejidad de esta función es \textbf{media}, porque usa dos
registros para la información de las asignaturas y los profesores de este
módulo y el número de tipo de datos elementales se encuentra entre seis y
diecinueve.

\subsubsection{Consultar relación asignaturas/aulas}
La complejidad de esta función es \textbf{media}, porque usa dos
registros para la información de las asignaturas y las aulas de este
módulo y el número de tipo de datos elementales se encuentra entre seis y
diecinueve.

\subsubsection{Consultar cuadrante de horarios}
La complejidad de esta función es \textbf{alta}, porque usa todos los
registros disponibles en este módulo y el número de tipo de datos elementales
es mayor que veinte.

\subsection{Archivos lógicos internos}
\subsubsection{Datos de las asignaturas}
La complejidad de esta función es \textbf{baja}, porque usa un
registros para la información de las asignaturas de este módulo y el número
de tipo de datos elementales es menor que veinte.

\subsubsection{Datos de los cursos}
La complejidad de esta función es \textbf{baja}, porque usa un
registros para la información de los cursos de este módulo y el número
de tipo de datos elementales es menor que veinte.

\subsubsection{Datos de las aulas}
La complejidad de esta función es \textbf{baja}, porque usa un
registros para la información de las aulas de este módulo y el número
de tipo de datos elementales es menor que veinte.

\subsection{Archivos lógicos de interfaz externos}
\subsubsection{Historial de guías docentes}
La complejidad de esta función es \textbf{alta}, porque usa todos los
registros de este módulo y el número de tipo de datos elementales es mayor
que cincuenta y uno. 

\section{Interfaz}
\subsection{Entradas externas}
\subsubsection{Autentificación de los usuarios}
La complejidad de esta función de usuario es \textbf{baja}, porque tiene un
registro para la información del usuario y el número de tipo de datos
elementales es dos. 

\subsubsection{Introducción de datos para la inserción/modificación/borrado 
  de información de un determinado usuario}
La complejidad para cada función de usuario (inserción, modificación y
borrado) es \textbf{media}, porque tiene un registro para la información del
usuario y el número de tipo de datos elementales es diecisiete.

\subsubsection{Introducción de datos para la inserción/modificación/borrado 
  de información de un determinado curso }
La complejidad para cada función de usuario (inserción, modificación y
borrado) es \textbf{baja}, porque tiene un registro para la información del
curso y el número de tipo de datos elementales es diez.

\subsection{Salidas externas}
\subsubsection{Mostrar los datos para una búsqueda de información con
  posibilidad de ordenación por diferentes criterios}
La complejidad de esta función de usuario es \textbf{media}, porque tiene un
registro para la información y el número de tipo de datos elementales es veinte.

\subsubsection{Mostrar los datos de un determinado usuario}
La complejidad de esta función de usuario es \textbf{baja}, porque tiene un
registro para la información del usuario y el número de tipo de datos
elementales es nueve.

\subsubsection{Mostrar los datos de un determinado curso}
La complejidad de esta función de usuario es \textbf{baja}, porque tiene un
registro para la información del curso y el número de tipo de datos
elementales es cinco.

\subsection{Consultas externas}
\subsubsection{Búsqueda de información con posibilidad de ordenación por
  diferentes criterios}
La complejidad de esta función de usuario es \textbf{media}, porque no tiene
ningún registro para la información y el número de tipo de datos elementales
es dieciocho.

\subsubsection{Consulta de los datos de un determinado usuario}
La complejidad de esta función de usuario es \textbf{baja}, porque no tiene
ningún para la información del usuario y el número de tipo de datos
elementales es uno.

\subsubsection{Consulta de los datos de un determinado curso}
La complejidad de esta función de usuario es \textbf{baja}, porque no tiene
ningún registro para la información del curso y el número de tipo de datos
elementales es uno.

\subsection{Archivos lógicos internos}
\subsubsection{Importación de datos desde un fichero}
La complejidad de esta función de usuario es \textbf{baja}, porque tiene un
fichero referenciado para la información y el número de tipo de datos
elementales es nueve.

\subsubsection{Exportación de los resultados de la búsqueda a un fichero}
La complejidad de esta función de usuario es \textbf{alta}, porque tiene dos
fichero referenciado para la información y el número de tipo de datos
elementales es veinte.

\chapter{Puntos función sin ajustar} \label{cap:pun}
Para la estimación del tamaño del producto software se necesitan los puntos
función sin ajustar representados en las tablas \ref{Tab:PFSAusu},
\ref{Tab:PFSAcur} y \ref{Tab:PFSAint}.

\section{Gestión de usuarios}
\begin{table}[!h]
  \centering
  \begin{tabular}{|b{2.1cm}<\centering|b{2.5cm}<{\centering}  ||c c c c c|}
    \hline
    \textbf{Tipo de función} & \textbf{Nivel de complejidad} & \textbf{Número}
    & \textbf{*} & \textbf{Peso} & \textbf{=} & \textbf{Total} \\
    \hline \hline
    \multirow{3}{*}{Entradas} 
    & Baja  & 1 & * & 3 & = & \textbf{3} \\ \cline{2-7}
    & Media & 1 & * & 4 & = & \textbf{4} \\ \cline{2-7}
    & Alta  &   &   & 6  &   & \\ \hline
    
    \multirow{3}{*}{Salidas}
    & Baja  & 1 & * & 4 & = & \textbf{4} \\ \cline{2-7}
    & Media & 3 & * & 5 & = & \textbf{15} \\ \cline{2-7}
    & Alta  & & & 7 & & \\ \hline

    \multirow{3}{*}{Consultas}
    & Baja  & 3 & * & 3 & = & \textbf{9} \\ \cline{2-7}
    & Media & & & 4 & & \\ \cline{2-7}
    & Alta  & & & 6 & & \\ \hline

    \multirow{3}{*}{Interfaces}
    & Baja  & 2 & * & 5 & = & \textbf{10}\\ \cline{2-7}
    & Media & & & 7 & & \\ \cline{2-7}
    & Alta  & & & 10 & & \\
    \hline \hline
    \multicolumn{6}{|l}{\textbf{Número de puntos función sin ajustar}} &
    \textbf{\textcolor{rojo}{45}} \\ 
    \hline
  \end{tabular}
  \caption{Puntos de función sin ajustar del módulo gestión de usuarios} 
  \label{Tab:PFSAusu}
\end{table}

\section{Gestión de cursos}
\begin{table}[!h]
  \centering
  \begin{tabular}{|b{2.1cm}<\centering|b{2.5cm}<{\centering} ||c c c c c|}
    \hline
    \textbf{Tipo de función} & \textbf{Nivel de complejidad} &
    \textbf{Número} & \textbf{*} & \textbf{Peso} & \textbf{=} & \textbf{Total}\\
    \hline \hline
    \multirow{3}{*}{Entradas} 
    & Baja & 4 & * & 3 & = & \textbf{12} \\ \cline{2-7}
    & Media &  &  & 4 &  &  \\ \cline{2-7}
    & Alta  &  &  & 6 & &  \\ \hline

    \multirow{3}{*}{Salidas}
    & Baja  & & & 4 & & \\ \cline{2-7}
    & Media & 3 & * & 5 & = & \textbf{15} \\ \cline{2-7}
    & Alta  & 1 & * & 7 & = & \textbf{7} \\ \hline

    \multirow{3}{*}{Consultas}
    & Baja  & & & 3 & & \\ \cline{2-7}
    & Media & 3 & * & 4 & = & \textbf{12} \\ \cline{2-7}
    & Alta  & 1 & * & 6 & = & \textbf{6} \\ \hline

    \multirow{3}{*}{Archivos}
    & Baja  & 3 & * & 7 & = & \textbf{21} \\ \cline{2-7}
    & Media &  &  & 10 &  & \\ \cline{2-7}
    & Alta  &  &  & 15 &  & \\ \hline

    \multirow{3}{*}{Interfaces}
    & Baja  & & & 5 & & \\ \cline{2-7}
    & Media & & & 7 & & \\ \cline{2-7}
    & Alta  & 1 & * & 10 & = & \textbf{10} \\
    \hline \hline
    \multicolumn{6}{|l}{\textbf{Número de puntos función sin ajustar}} &
    \textbf{\textcolor{rojo}{83}} \\ 
    \hline
  \end{tabular}
  \caption{Puntos de función sin ajustar del módulo gestión de cursos} 
  \label{Tab:PFSAcur}
\end{table}

\section{Interfaz}
\begin{table}[!h]
  \centering
  \begin{tabular}{|b{2.1cm}<\centering|b{2.5cm}<{\centering} ||c c c c c|}
    \hline
    \textbf{Tipo de función} & \textbf{Nivel de complejidad} &
    \textbf{Número} & \textbf{*} & \textbf{Peso} & \textbf{=} & \textbf{Total}\\
        \hline \hline
    \multirow{3}{*}{Entradas} 
    & Baja & 4 & * & 3 & = & \textbf{12} \\ \cline{2-7}
    & Media & 3 & * & 4 & = & \textbf{12} \\ \cline{2-7}
    & Alta  & & & 6 & & \\ \hline

    \multirow{3}{*}{Salidas}
    & Baja  & 2 & * & 4 & = & \textbf{8} \\ \cline{2-7}
    & Media & 1 & * & 5 & = & \textbf{5} \\ \cline{2-7}
    & Alta  & & & 7 & & \\ \hline

    \multirow{3}{*}{Consultas}
    & Baja  & 2 & * & 3 & = & \textbf{6} \\ \cline{2-7}
    & Media & 1 & * & 4 & = & \textbf{4} \\ \cline{2-7}
    & Alta  & & & 6 & & \\ \hline \hline

    \multirow{3}{*}{Archivos}
    & Baja  & 1 & * & 7 & = & \textbf{7} \\ \cline{2-7}
    & Media & & & 10 & & \\ \cline{2-7}
    & Alta  & 1 & * & 15 & = & \textbf{15} \\
    \hline \hline
    \multicolumn{6}{|l}{\textbf{Número de puntos función sin ajustar}} &
    \textbf{\textcolor{rojo}{69}} \\ 
    \hline
  \end{tabular}
  \caption{Puntos de función sin ajustar del módulo interfaz} 
  \label{Tab:PFSAint}
\end{table}

\chapter{Factores de influencia} \label{cap:fac}
En este apartado se especifican los pesos que se han asignado a cada factor
de influencia. Estos valores han sido extraídos de la información que
proporciona Métrica V3 \cite{met}. 

\section{Gestión de usuarios}
\begin{enumerate}[{\bf 1.}]

\item {\bf Comunicación de datos:} \emph{Peso 4} \\ La aplicación está basada en
  un teleproceso interactivo, pero con un solo protocolo de comunicaciones.

\item {\bf Procesamiento distribuido:} \emph{Peso 4} \\ Proceso distribuido
  con transferencia de datos on line en ambas direcciones. 

\item {\bf Objetivos de rendimiento:} \emph{Peso 1} \\ Se establecen requisitos
  de rendimiento y diseño, pero no se requieren acciones especiales.

\item {\bf Configuración de uso intensivo:} \emph{Peso 2} \\ Se incluyen algunas
  consideraciones sobre tiempo y seguridad. 

\item {\bf Tasas de transacción rápida:} \emph{Peso 0} \\ Las transacciones no
  están afectadas por picos de tráfico. 

\item {\bf Entradas de datos en línea:} \emph{Peso 5} \\ Más del 30\% de las
  funciones son entradas interactivas de datos. 

\item {\bf Eficiencia de Usuario Final:} \emph{Peso 5} \\ Además se requiere del
  uso de herramientas y procesos para demostrar que se cumplen los requisitos
  establecidos.

\item {\bf Actualización de datos en línea:} \emph{Peso 5} \\ Además, grandes
  volúmenes implican consideraciones de coste en el proceso de
  recuperación. Se incluyen procedimientos de recuperación que requieren
  mínima intervención del operador. 

\item {\bf Procesamiento complejo:} \emph{Peso 1} \\ Es aplicable una de las
  situaciones descritas. 

\item {\bf Reutilización:} \emph{Peso 3} \\ El 10\% o más de la aplicación se
  considera reusable. 

\item {\bf Facilidad de instalación:} \emph{Peso 2} \\ El usuario ha declarado
  consideraciones especiales para la conversión e instalación y se requieren
  guías probadas de conversión e instalación. El impacto de la conversión
  sobre el proyecto no se considera importante. 

\item {\bf Facilidad operacional:} \emph{Peso 5} \\ La aplicación debe diseñarse
  para una operación totalmente automática. La recuperación automática ante
  errores es una funcionalidad de la aplicación. 

\item {\bf Múltiples localizaciones:} \emph{Peso 2} \\ Se consideran en el
  diseño múltiples instalaciones pero para operar con similar configuración HW
  y/o SW.

\item {\bf Facilidad de cambio:} \emph{Peso 1} Consultas y Generación de
  Informes Flexible que manejan peticiones simples.

\end{enumerate}

Todo lo anterior queda resumido en la tabla \ref{tab:infUS}

\begin{table}[!ht]
  \centering
  \begin{tabular}{|l|c|}
    \hline
    \textbf{Factores de influencia en la dificultad del sistema} & 
    \textbf{GI} \\
    \hline \hline
    Comunicación de datos           & 4 \\ \hline
    Procesamiento distribuido       & 4 \\ \hline
    Objetivos de rendimiento        & 1 \\ \hline
    Configuración de uso intensivo  & 2 \\ \hline
    Tasas de transacción rápida     & 0 \\ \hline
    Entradas de datos en línea      & 5 \\ \hline
    Eficiencia de Usuario Final     & 5 \\ \hline
    Actualización de datos en línea & 5 \\ \hline
    Procesamiento complejo          & 1 \\ \hline
    Reutilización                   & 3 \\ \hline
    Facilidad de instalación        & 2 \\ \hline
    Facilidad operacional           & 5 \\ \hline
    Múltiples localizaciones        & 2 \\ \hline
    Facilidad de cambio             & 1 \\ \hline \hline
    \textbf{Total}                  & \textbf{40} \\ \hline
  \end{tabular}
  \label{tab:infUS}
  \caption{Influencia en la dificultad del sistema del módulo usuarios}
\end{table}


\section{Gestión de cursos}
\begin{enumerate}[{\bf 1.}]

\item {\bf Comunicación de datos:} \emph{Peso 4} \\ La aplicación está basada en
  un teleproceso interactivo, pero con un solo protocolo de comunicaciones.

\item {\bf Procesamiento distribuido:} \emph{Peso 4} \\ Proceso distribuido
  con transferencia de datos on line en ambas direcciones. 

\item {\bf Objetivos de rendimiento:} \emph{Peso 1} \\ Se establecen requisitos
  de rendimiento y diseño, pero no se requieren acciones especiales.

\item {\bf Configuración de uso intensivo:} \emph{Peso 2} \\ Se incluyen algunas
  consideraciones sobre tiempo y seguridad. 

\item {\bf Tasas de transacción rápida:} \emph{Peso 0} \\ Las transacciones no
  están afectadas por picos de tráfico. 

\item {\bf Entradas de datos en línea:} \emph{Peso 5} \\ Más del 30\% de las
  funciones son entradas interactivas de datos. 

\item {\bf Eficiencia de Usuario Final:} \emph{Peso 5} \\ Además se requiere del
  uso de herramientas y procesos para demostrar que se cumplen los requisitos
  establecidos.

\item {\bf Actualización de datos en línea:} \emph{Peso 5} \\ Además, grandes
  volúmenes implican consideraciones de coste en el proceso de
  recuperación. Se incluyen procedimientos de recuperación que requieren
  mínima intervención del operador. 

\item {\bf Procesamiento complejo:} \emph{Peso 5} \\Son aplicables todas las
  situaciones descritas.
 
\item {\bf Reutilización:} \emph{Peso 3} \\ El 10\% o más de la aplicación se
  considera reusable. 

\item {\bf Facilidad de instalación:} \emph{Peso 3} \\ El usuario ha declarado
  consideraciones especiales para la conversión e instalación y se requieren
  guías probadas de conversión e instalación. El impacto de la conversión
  sobre el proyecto se considera importante. 

\item {\bf Facilidad operacional:} \emph{Peso 5} \\ La aplicación debe diseñarse
  para una operación totalmente automática. La recuperación automática ante
  errores es una funcionalidad de la aplicación. 

\item {\bf Múltiples localizaciones:} \emph{Peso 2} \\ Se consideran en el
  diseño múltiples instalaciones pero para operar con similar configuración HW
  y/o SW.

\item {\bf Facilidad de cambio:} \emph{Peso 3} \\ Consultas y Generación de
  Informes Flexible que manejan peticiones complejas.

Todo lo anterior queda resumido en la tabla \ref{tab:infCU}

\end{enumerate}
\begin{table}[!ht]
  \centering
  \begin{tabular}{|l|c|}
    \hline
    \textbf{Factores de influencia en la dificultad del sistema} & 
    \textbf{GI} \\
    \hline \hline
    Comunicación de datos           & 4 \\ \hline
    Procesamiento distribuido       & 4 \\ \hline
    Objetivos de rendimiento        & 1 \\ \hline
    Configuración de uso intensivo  & 2 \\ \hline
    Tasas de transacción rápida     & 0 \\ \hline
    Entradas de datos en línea      & 5 \\ \hline
    Eficiencia de Usuario Final     & 5 \\ \hline
    Actualización de datos en línea & 5 \\ \hline
    Procesamiento complejo          & 5 \\ \hline
    Reutilización                   & 3 \\ \hline
    Facilidad de instalación        & 3 \\ \hline
    Facilidad operacional           & 5 \\ \hline
    Múltiples localizaciones        & 2 \\ \hline
    Facilidad de cambio             & 3 \\ \hline \hline
    \textbf{Total}                  & \textbf{47} \\ \hline
  \end{tabular}
  \label{tab:infCU}
  \caption{Influencia en la dificultad del sistema del módulo cursos}
\end{table}

\section{Interfaz}
\begin{enumerate}[{\bf 1.}]

\item {\bf Comunicación de datos:} \emph{Peso 4} \\ La aplicación está basada en
  un teleproceso interactivo, pero con un solo protocolo de comunicaciones.

\item {\bf Procesamiento distribuido:} \emph{Peso 4} \\ Proceso distribuido
  con transferencia de datos on line en ambas direcciones. 

\item {\bf Objetivos de rendimiento:} \emph{Peso 1} \\ Se establecen requisitos
  de rendimiento y diseño, pero no se requieren acciones especiales.

\item {\bf Configuración de uso intensivo:} \emph{Peso 2} \\ Se incluyen algunas
  consideraciones sobre tiempo y seguridad. 

\item {\bf Tasas de transacción rápida:} \emph{Peso 0} \\ Las transacciones no
  están afectadas por picos de tráfico. 

\item {\bf Entradas de datos en línea:} \emph{Peso 5} \\ Más del 30\% de las
  funciones son entradas interactivas de datos. 

\item {\bf Eficiencia de Usuario Final:} \emph{Peso 5} \\ Además se requiere del
  uso de herramientas y procesos para demostrar que se cumplen los requisitos
  establecidos.

\item {\bf Actualización de datos en línea:} \emph{Peso 5} \\ Además, grandes
  volúmenes implican consideraciones de coste en el proceso de
  recuperación. Se incluyen procedimientos de recuperación que requieren
  mínima intervención del operador. 

\item {\bf Procesamiento complejo:} \emph{Peso 1} \\ Es aplicable una de las
  situaciones descritas. 

\item {\bf Reutilización:} \emph{Peso 3} \\ El 10\% o más de la aplicación se
  considera reusable. 

\item {\bf Facilidad de instalación:} \emph{Peso 1} \\ El usuario no ha
  declarado consideraciones especiales pero se requiere una configuración
  especial para la instalación 

\item {\bf Facilidad operacional:} \emph{Peso 5} \\ La aplicación debe diseñarse
  para una operación totalmente automática. La recuperación automática ante
  errores es una funcionalidad de la aplicación. 

\item {\bf Múltiples localizaciones:} \emph{Peso 2} \\ Se consideran en el
  diseño múltiples instalaciones pero para operar con similar configuración
  HW y/o SW. 

\item {\bf Facilidad de cambio:} \emph{Peso 1} \\ Consultas y Generación de
  Informes Flexible que manejan peticiones simples.
\end{enumerate}

Todo lo anterior queda resumido en la tabla \ref{tab:infIN}

\begin{table}[!ht]
  \centering
  \begin{tabular}{|l|c|}
    \hline
    \textbf{Factores de influencia en la dificultad del sistema} & 
    \textbf{Peso} \\
    \hline \hline
    Comunicación de datos           & 4 \\ \hline
    Procesamiento distribuido       & 4 \\ \hline
    Objetivos de rendimiento        & 1 \\ \hline
    Configuración de uso intensivo  & 2 \\ \hline
    Tasas de transacción rápida     & 0 \\ \hline
    Entradas de datos en línea      & 5 \\ \hline
    Eficiencia de Usuario Final     & 5 \\ \hline
    Actualización de datos en línea & 5 \\ \hline
    Procesamiento complejo          & 1 \\ \hline
    Reutilización                   & 3 \\ \hline
    Facilidad de instalación        & 1 \\ \hline
    Facilidad operacional           & 5 \\ \hline
    Múltiples localizaciones        & 2 \\ \hline
    Facilidad de cambio             & 1 \\ \hline \hline
    \textbf{Total}                  & \textbf{39} \\ \hline
  \end{tabular}
  \label{tab:infIN}
  \caption{Influencia en la dificultad del sistema del módulo interfaz}
\end{table}

\chapter{Puntos función ajustados} \label{cap:aju}
Para el cálculo de los puntos función totales se utiliza la fórmula \ref{eq},
obteniendo como resultado los valores reflejados en la tabla \ref{Tab:PFA}.
\begin{equation}
  PF=PFSA*(0.65+0.01*SUMA(GI)
  \label{eq}
\end{equation}
\begin{table}[!h]
  \centering
  \begin{tabular}{|b{3.3cm}|b{2.9cm}<\centering|c|b{2.9cm}<\centering|}
    \hline
    \center{\textbf{Módulo}} & \textbf{Puntos función sin ajustar} &
    \textbf{Total GI} & \textbf{Puntos función ajustados}\\
    \hline \hline
    Gestión de usuarios & 45 & 40 & 47.25 \\ \hline
    Gestión de cursos   & 83 & 47 & 92.96 \\ \hline
    Interfaz            & 69 & 39 & 71.76 \\ \hline
  \end{tabular}
  \caption{Puntos de función ajustados} 
  \label{Tab:PFA}
\end{table}

\bibliographystyle{plain} 
\bibliography{t5}

\end{document}
